
\begin{document}

\setlength{\headheight}{15pt}
\vspace*{-2.5cm}
\begin{figure}[htbp]
  \begin{minipage}[r][][t]{0.15\textwidth}
\begin{center}
 \includegraphics[width=2.2cm, height=2.2cm]{imagen/inglogo.jpg}
\vspace{1.8cm}
\end{center}
\end{minipage}\quad
\begin{minipage}[r][][t]{0.31\textwidth}
\begin{center}
\vspace*{-25mm}
{\bf \materia }\\[2mm]
{\bf Módulo 4}
\end{center}
\end{minipage}\quad\begin{minipage}[r][][t]{0.31\textwidth}
\begin{center}
\vspace*{-25mm}
\textit{Ingeniería en Petróleo} \\[2mm]
Recursado 2026
\end{center}
\end{minipage}\quad
\begin{minipage}[r][][t]{0.15\textwidth}
  \begin{center}
     \includegraphics[width=2.2cm, height=2.2cm]{imagen/Logo_UNco.jpg}
     \vspace{1.8cm}
  \end{center}
\end{minipage}
\end{figure}
\vspace*{-33mm}

\begin{center}
{\bf \Large{ Resolución Trabajo Práctico 4: Vectores. Recursado }}
\end{center}

\newcommand{\ejesXY}[4]{%
\draw[->] (#1-.3,0) -- (#2+.3,0) node[right] {$x$};
\draw[->] (0,#3-.3) -- (0,#4+.3) node[above] {$y$};
\draw[color=gray!30,dash pattern=on 1pt off 1pt,xstep=1cm,ystep=1cm] (#1-.3,#3-.3) grid (#2+.3,#4+.3);
\foreach \x in {#1,...,-1,1,...,#2}{\draw[shift={(\x,0)},color=black] (0pt,2pt)--(0pt,-2pt) node[below] {\tiny $\x$};}
\foreach \y in {#3,...,-1,1,...,#4}{\draw[shift={(0,\y)},color=black] (2pt,0pt)--(-2pt,0pt) node[left] {\tiny $\y$};}
}
\newcommand{\dibvec}[4][blue]{\draw[->,line width=1.4pt,color=#1] #2 -- #3 node[pos=.62,above,sloped] {\scriptsize #4};}

\begin{enumerate}[\hspace{-3mm}$1)$]

\item Dados $P=(2,1)$, $Q=(6,3)$, $A=(-1,2)$ y $B=(5,-1)$.

\begin{sol}
El vector dado tiene componentes
\[
\overrightarrow{PQ}=Q-P=(6,3)-(2,1)=(4,2).
\]
Por lo tanto:
\begin{itemize}
\item[a)] Un representante con origen en $(0,0)$ es $\vec v=(4,2)$.
\item[b)] El vector con origen en $A$ e igual dirección y doble longitud es $2\vec v=(8,4)$. Su extremo es
\[
A+2\vec v=(-1,2)+(8,4)=(7,6).
\]
No es único si sólo se pide igual dirección y doble longitud, porque también puede tomarse el sentido opuesto $-2\vec v=(-8,-4)$. Si se exige además el mismo sentido, sí queda determinado.
\item[c)] El vector pedido es $\vec z=-\dfrac12\vec v=(-2,-1)$. Su extremo es
\[
B+\vec z=(5,-1)+(-2,-1)=(3,-2).
\]
\item[d)] El simétrico respecto del eje $x$ es $\vec v_x=(4,-2)$ y el simétrico respecto del eje $y$ es $\vec v_y=(-4,2)$.
\end{itemize}

\begin{center}
\begin{tikzpicture}[line cap=round,line join=round,>=triangle 45,scale=.62]
\ejesXY{-5}{9}{-4}{7}
\coordinate (P) at (2,1);\coordinate (Q) at (6,3);\coordinate (A) at (-1,2);\coordinate (B) at (5,-1);
\draw[fill=black] (P) circle (1.5pt) node[below right]{\scriptsize $P$};
\draw[fill=black] (Q) circle (1.5pt) node[above right]{\scriptsize $Q$};
\draw[fill=black] (A) circle (1.5pt) node[above left]{\scriptsize $A$};
\draw[fill=black] (B) circle (1.5pt) node[below right]{\scriptsize $B$};
\dibvec[blue]{(P)}{(Q)}{$\overrightarrow{PQ}$}
\dibvec[green!60!black]{(0,0)}{(4,2)}{$\vec v$}
\dibvec[red]{(A)}{(7,6)}{$2\vec v$}
\dibvec[teal]{(B)}{(3,-2)}{$\vec z$}
\dibvec[orange]{(0,0)}{(4,-2)}{$\vec v_x$}
\dibvec[purple]{(0,0)}{(-4,2)}{$\vec v_y$}
\end{tikzpicture}
\end{center}
\end{sol}

\item Dados $\vec u=(5,3)$, $\vec v=(-2,4)$ y $\vec w=(-4,-1)$.

\begin{sol}
Primero los representamos desde el origen:
\begin{center}
\begin{tikzpicture}[>=triangle 45,scale=.55]
\ejesXY{-9}{6}{-6}{17}
\dibvec[blue]{(0,0)}{(5,3)}{$\vec u$}
\dibvec[red]{(0,0)}{(-2,4)}{$\vec v$}
\dibvec[green!60!black]{(0,0)}{(-4,-1)}{$\vec w$}
\end{tikzpicture}
\end{center}

\begin{enumerate}[i)]
\item $\vec w+2\vec v=(-4,-1)+2(-2,4)=(-4,-1)+(-4,8)=(-8,7)$.
\item $-\dfrac12\vec u+\vec v=-\dfrac12(5,3)+(-2,4)=\left(-\dfrac52,-\dfrac32\right)+(-2,4)=\left(-\dfrac92,\dfrac52\right)$.
\item $\dfrac12(\vec v+\vec u)=\dfrac12\big((-2,4)+(5,3)\big)=\dfrac12(3,7)=\left(\dfrac32,\dfrac72\right)$.
\item $-\vec v+\vec w=(2,-4)+(-4,-1)=(-2,-5)$.
\item $\dfrac12\vec v+\dfrac12\vec u=(-1,2)+\left(\dfrac52,\dfrac32\right)=\left(\dfrac32,\dfrac72\right)$.
\item $\vec u+3\vec v-\vec w=(5,3)+3(-2,4)-(-4,-1)=(3,16)$.
\end{enumerate}

Los incisos III) y V) se relacionan por la propiedad distributiva del producto de un escalar respecto de la suma:
\[
\frac12(\vec v+\vec u)=\frac12\vec v+\frac12\vec u.
\]

\begin{center}
\begin{tikzpicture}[>=triangle 45,scale=.48]
\ejesXY{-9}{6}{-6}{17}
\dibvec[gray]{(0,0)}{(5,3)}{$\vec u$}
\dibvec[gray]{(0,0)}{(-2,4)}{$\vec v$}
\dibvec[gray]{(0,0)}{(-4,-1)}{$\vec w$}
\dibvec[red]{(0,0)}{(-8,7)}{$\vec w+2\vec v$}
\dibvec[blue]{(0,0)}{(-4.5,2.5)}{$-\frac12\vec u+\vec v$}
\dibvec[green!60!black]{(0,0)}{(1.5,3.5)}{$\frac12(\vec v+\vec u)$}
\dibvec[orange]{(0,0)}{(-2,-5)}{$-\vec v+\vec w$}
\dibvec[purple]{(0,0)}{(3,16)}{$\vec u+3\vec v-\vec w$}
\end{tikzpicture}
\end{center}
\end{sol}

\item Dados $\vec v=(3,-6,2)$, $\vec w=\left(\dfrac15,0,-2\right)$ y $\vec z=(-2,4,1)$.
\begin{sol}
\begin{enumerate}[a)]
\item
\begin{enumerate}[i)]
\item
\[
2(\vec v-3\vec w)=2\left((3,-6,2)-3\left(\frac15,0,-2\right)\right)
=2\left(\frac{12}{5},-6,8\right)=\left(\frac{24}{5},-12,16\right).
\]
\item
\[
\begin{aligned}
-3\left(\frac12\vec z+\vec v\right)+2\vec w
&=-3\left((-1,2,\frac12)+(3,-6,2)\right)+\left(\frac25,0,-4\right)\\
&=-3\left(2,-4,\frac52\right)+\left(\frac25,0,-4\right)\\
&=(-6,12,-\frac{15}{2})+\left(\frac25,0,-4\right)\\
&=\left(-\frac{28}{5},12,-\frac{23}{2}\right).
\end{aligned}
\]
\end{enumerate}
\item Despejamos $\vec u$:
\[
\vec v-2\vec u+\vec z=2\vec w-4\vec u
\quad\Rightarrow\quad
2\vec u=2\vec w-\vec v-\vec z.
\]
Entonces,
\[
\vec u=\vec w-\frac12\vec v-\frac12\vec z
=\left(\frac15,0,-2\right)-\left(\frac32,-3,1\right)-(-1,2,\frac12)
=\left(-\frac{3}{10},1,-\frac72\right).
\]
\end{enumerate}
\end{sol}

\item Los siete vectores están sobre una circunferencia y tienen el mismo módulo.
\begin{sol}
Los pares opuestos se anulan:
\[
\vec t+\vec p=\vec 0,\qquad \vec m+\vec q=\vec 0,\qquad \vec u+\vec r=\vec 0.
\]
Queda solamente el vector $\vec n$. Por lo tanto, la respuesta correcta es:
\[
\boxed{\text{c) Es igual a }\vec n.}
\]
\begin{center}
\begin{tikzpicture}[>=triangle 45,rotate=-15]
\xdef\r{2.1}
\draw[dashed] (0,0) circle(\r cm);
\foreach \angle/\label in {0/$\vec t$,45/$\vec m$,90/$\vec n$,135/$\vec u$,180/$\vec p$,225/$\vec q$,315/$\vec r$}{
\draw[->,thick] (0,0)--(\angle:\r cm) node[pos=1.15]{\label};}
\draw[->,line width=2pt,blue] (0,0)--(90:\r cm) node[pos=.55,right]{$\vec n$};
\fill[blue] (0,0) circle(2pt);
\end{tikzpicture}
\end{center}
\end{sol}

\item Determinar componentes de $\overrightarrow{PQ}$.
\begin{sol}
\begin{enumerate}[a)]
\item $\overrightarrow{PQ}=Q-P=(2,-1)-(-3,4)=(5,-5)$.
\item $\overrightarrow{PQ}=Q-P=(3,-2,1)-(-1,2,4)=(4,-4,-3)$.
\end{enumerate}
\end{sol}

\item Dados $A=(2,1,-1)$, $B=(3,-2,0)$ y $C=(-1,4,2)$.
\begin{sol}
\begin{enumerate}[a)]
\item
\[
\overrightarrow{BC}=C-B=(-1,4,2)-(3,-2,0)=(-4,6,2).
\]
Como $\overrightarrow{AD}=\overrightarrow{BC}$, entonces
\[
D=A+\overrightarrow{BC}=(2,1,-1)+(-4,6,2)=(-2,7,1).
\]
\item
\[
\overrightarrow{AB}=B-A=(3,-2,0)-(2,1,-1)=(1,-3,1).
\]
Para que $\overrightarrow{CM}\parallel \overrightarrow{AB}$ debe existir $k\in\mathbb R$, $k\neq0$, tal que
\[
\overrightarrow{CM}=k\overrightarrow{AB}.
\]
Por lo tanto,
\[
M=C+k(1,-3,1)=(-1+k,4-3k,2+k),\qquad k\in\mathbb R,\ k\neq0.
\]
No es único, porque hay infinitos valores posibles de $k$.
\end{enumerate}
\end{sol}

\item Normas.
\begin{sol}
\begin{enumerate}[a)]
\item $\|\vec a\|=\sqrt{2^2+\left(\frac32\right)^2}=\sqrt{\frac{25}{4}}=\dfrac52$.
\item $\|-\vec a\|=\|\vec a\|=\dfrac52$.
\item $\vec b-\vec a=(-1,5)-\left(2,\frac32\right)=\left(-3,\frac72\right)$, entonces
\[
\|\vec b-\vec a\|=\sqrt{(-3)^2+\left(\frac72\right)^2}=\sqrt{\frac{85}{4}}=\frac{\sqrt{85}}{2}.
\]
\item $\|\vec c\|=\sqrt{(2\sqrt2)^2+(-\sqrt2)^2+(\sqrt2)^2}=\sqrt{12}=2\sqrt3$. Luego,
\[
\|-2\vec c\|=2\|\vec c\|=4\sqrt3.
\]
\end{enumerate}
\end{sol}

\item Valores de $k$.
\begin{sol}
\begin{enumerate}[a)]
\item
\[
\|(2,k,1)\|=3\Rightarrow \sqrt{4+k^2+1}=3\Rightarrow k^2+5=9\Rightarrow k^2=4.
\]
Entonces $\boxed{k=2\text{ o }k=-2}$.
\item
\[
\|k(1,2,2)\|=|k|\sqrt{1+4+4}=3|k|.
\]
Como $3|k|=6$, resulta $|k|=2$. Entonces $\boxed{k=2\text{ o }k=-2}$.
\end{enumerate}
\end{sol}

\item Distancias.
\begin{sol}
\begin{enumerate}[a)]
\item
\[
d(P,Q)=\sqrt{(4-(-2))^2+(-2-1)^2}=\sqrt{36+9}=3\sqrt5.
\]
\item
\[
d(P,Q)=\sqrt{(-1-2)^2+(3-(-1))^2+(2-0)^2}=\sqrt{9+16+4}=\sqrt{29}.
\]
\end{enumerate}
\end{sol}

\item Sean $\vec u=(2,-1,2)$ y $\vec p=(2,1,2)$.
\begin{sol}
\begin{enumerate}[a)]
\item Si $\vec w\parallel \vec p$, entonces $\vec w=k\vec p$. Como $\|\vec p\|=\sqrt{4+1+4}=3$, se tiene
\[
\|k\vec p\|=|k|3=6\Rightarrow |k|=2.
\]
Los vectores son
\[
\vec w_1=2\vec p=(4,2,4),\qquad \vec w_2=-2\vec p=(-4,-2,-4).
\]
Existen exactamente dos.
\item El opuesto de $\dfrac12\vec u$ es
\[
-\frac12\vec u=-\frac12(2,-1,2)=\left(-1,\frac12,-1\right).
\]
\item
\[
\vec u+\vec p=(2,-1,2)+(2,1,2)=(4,0,4).
\]
Su norma es $\sqrt{16+16}=4\sqrt2$. El versor de sentido contrario es
\[
-\frac{\vec u+\vec p}{\|\vec u+\vec p\|}=-\frac{(4,0,4)}{4\sqrt2}=\left(-\frac1{\sqrt2},0,-\frac1{\sqrt2}\right).
\]
\end{enumerate}
\end{sol}

\item Versores canónicos.
\begin{sol}
\begin{enumerate}[a)]
\item En $\mathbb R^2$:
\[
\vec u=\left(3,-\frac52\right)=3\iv-\frac52\jv.
\]
\item En $\mathbb R^3$:
\[
\vec v=(2,-1,6)=2\iv-\jv+6\kv.
\]
\end{enumerate}
\end{sol}

\item Componentes a partir de módulo y ángulo.
\begin{sol}
Usamos $\vec u=\|\vec u\|(\cos\theta,\sin\theta)$.
\begin{enumerate}[a)]
\item Para $\|\vec u\|=8$ y $\theta=30^\circ$:
\[
\vec u=8\left(\cos30^\circ,\sin30^\circ\right)=8\left(\frac{\sqrt3}{2},\frac12\right)=(4\sqrt3,4).
\]
\item Para $\|\vec u\|=6$ y $\theta=150^\circ$:
\[
\vec u=6\left(\cos150^\circ,\sin150^\circ\right)=6\left(-\frac{\sqrt3}{2},\frac12\right)=(-3\sqrt3,3).
\]
\end{enumerate}
\begin{center}
\begin{tikzpicture}[>=triangle 45,scale=.45]
\ejesXY{-6}{8}{-1}{6}
\dibvec[blue]{(0,0)}{(6.928,4)}{$(4\sqrt3,4)$}
\dibvec[red]{(0,0)}{(-5.196,3)}{$(-3\sqrt3,3)$}
\end{tikzpicture}
\end{center}
\end{sol}

\item Fluido hacia el noroeste.
\begin{sol}
La dirección noroeste forma un ángulo de $135^\circ$ medido desde el eje positivo $x$ (este). Entonces
\[
\vec v=8(\cos135^\circ,\sin135^\circ)
=8\left(-\frac{\sqrt2}{2},\frac{\sqrt2}{2}\right)=(-4\sqrt2,4\sqrt2).
\]
La componente hacia el este es $-4\sqrt2\,\text{m/s}$, es decir $4\sqrt2\,\text{m/s}$ hacia el oeste, y la componente hacia el norte es $4\sqrt2\,\text{m/s}$.
\end{sol}

\item Fuerza resultante en tuberías.
\begin{sol}
Tomamos este como eje $x$ positivo y norte como eje $y$ positivo. Las fuerzas son
\[
\vec F_1=(-400,0),\qquad \vec F_2=(0,300).
\]
La resultante es
\[
\vec R=\vec F_1+\vec F_2=(-400,300).
\]
Su módulo es
\[
\|\vec R\|=\sqrt{(-400)^2+300^2}=\sqrt{250000}=500\,\text{N}.
\]
La dirección cumple
\[
\tan\theta=\frac{300}{400}=\frac34\Rightarrow \theta\approx36,87^\circ.
\]
Por lo tanto, la fuerza resultante es de $500\,\text{N}$ con dirección $36,87^\circ$ al norte del oeste.
\end{sol}

\item Velocidad resultante del buque.
\begin{sol}
\[
\vec v_1=(12,0),\qquad \vec v_2=(0,5).
\]
Entonces
\[
\vec v_R=(12,5),\qquad \|\vec v_R\|=\sqrt{12^2+5^2}=13\,\text{m/s}.
\]
La dirección respecto del este es
\[
\theta=\arctan\left(\frac5{12}\right)\approx22,62^\circ.
\]
Por lo tanto, el buque se mueve a $13\,\text{m/s}$ con dirección $22,62^\circ$ al norte del este.
\end{sol}

\item Producto escalar.
\begin{sol}
\begin{enumerate}[a)]
\item
\[
\vec u\cdot\vec v=(4,-3)\cdot(-2,5)=4(-2)+(-3)5=-8-15=-23.
\]
\item
\[
\vec u\cdot\vec v=\|\vec u\|\,\|\vec v\|\cos\theta=6\cdot5\cdot\cos\left(\frac\pi6\right)=30\cdot\frac{\sqrt3}{2}=15\sqrt3.
\]
\end{enumerate}
\end{sol}

\item Dados $\vec a=(-3,\alpha)$ y $\vec b=(\beta,2)$.
\begin{sol}
De la condición $\|\vec a\|=5$:
\[
\sqrt{(-3)^2+\alpha^2}=5\Rightarrow 9+\alpha^2=25\Rightarrow \alpha^2=16.
\]
Luego $\alpha=4$ o $\alpha=-4$.
Además,
\[
\vec a\cdot\vec b=5\Rightarrow (-3,\alpha)\cdot(\beta,2)=5\Rightarrow -3\beta+2\alpha=5.
\]
Si $\alpha=4$:
\[
-3\beta+8=5\Rightarrow \beta=1.
\]
Si $\alpha=-4$:
\[
-3\beta-8=5\Rightarrow \beta=-\frac{13}{3}.
\]
Por lo tanto,
\[
\boxed{(\alpha,\beta)=(4,1)\quad\text{o}\quad (\alpha,\beta)=\left(-4,-\frac{13}{3}\right).}
\]
\end{sol}

\item Cubo de arista $2\,\text{m}$.
\begin{sol}
Podemos ubicar una diagonal principal como el vector $\vec d=(2,2,2)$ y una diagonal de la base como $\vec b=(2,2,0)$.
\begin{enumerate}[a)]
\item
\[
\|\vec d\|=\sqrt{2^2+2^2+2^2}=2\sqrt3\,\text{m}.
\]
\item
\[
\cos\theta=\frac{\vec d\cdot\vec b}{\|\vec d\|\,\|\vec b\|}
=\frac{(2,2,2)\cdot(2,2,0)}{(2\sqrt3)(2\sqrt2)}
=\frac{8}{4\sqrt6}=\frac{\sqrt6}{3}.
\]
Entonces
\[
\theta=\arccos\left(\frac{\sqrt6}{3}\right)\approx35,26^\circ.
\]
\end{enumerate}
\end{sol}

\item Ángulo y proyección $\operatorname{Proy}_{\vec s}\vec r$.
\begin{sol}
Recordamos que
\[
\cos\theta=\frac{\vec r\cdot\vec s}{\|\vec r\|\,\|\vec s\|},\qquad
\operatorname{Proy}_{\vec s}\vec r=\frac{\vec r\cdot\vec s}{\|\vec s\|^2}\vec s.
\]
\begin{enumerate}[a)]
\item $\vec r=(2,0)$ y $\vec s=(1,\sqrt3)$.
\[
\vec r\cdot\vec s=2,\quad \|\vec r\|=2,\quad \|\vec s\|=2.
\]
Entonces
\[
\cos\theta=\frac{2}{4}=\frac12\Rightarrow \theta=\frac\pi3.
\]
La proyección es
\[
\operatorname{Proy}_{\vec s}\vec r=\frac{2}{4}(1,\sqrt3)=\left(\frac12,\frac{\sqrt3}{2}\right).
\]
\begin{center}
\begin{tikzpicture}[>=triangle 45,scale=.9]
\ejesXY{-1}{3}{-1}{3}
\dibvec[blue]{(0,0)}{(2,0)}{$\vec r$}
\dibvec[red]{(0,0)}{(1,1.732)}{$\vec s$}
\dibvec[green!60!black]{(0,0)}{(.5,.866)}{$\operatorname{Proy}_{\vec s}\vec r$}
\draw[->,thick] (.7,0) arc[start angle=0,end angle=60,radius=.7cm] node[pos=.6,right]{\tiny $\pi/3$};
\end{tikzpicture}
\end{center}
\item $\vec r=(1,2,2)$ y $\vec s=(2,1,-2)$.
\[
\vec r\cdot\vec s=1\cdot2+2\cdot1+2(-2)=0.
\]
Como ambos son no nulos, forman un ángulo recto:
\[
\theta=\frac\pi2,\qquad \operatorname{Proy}_{\vec s}\vec r=(0,0,0).
\]
\item $\vec r=(2,-1,2)$ y $\vec s=(1,2,2)$.
\[
\vec r\cdot\vec s=2-2+4=4,
\quad \|\vec r\|=3,
\quad \|\vec s\|=3.
\]
Entonces
\[
\cos\theta=\frac49\Rightarrow \theta=\arccos\left(\frac49\right).
\]
La proyección es
\[
\operatorname{Proy}_{\vec s}\vec r=\frac{4}{9}(1,2,2)=\left(\frac49,\frac89,\frac89\right).
\]
\end{enumerate}
\end{sol}

\item Trabajo.
\begin{sol}
\begin{enumerate}[a)]
\item El desplazamiento es
\[
\vec r=B-A=(3,-2,1)-(1,0,-2)=(2,-2,3).
\]
Entonces
\[
W=\vec F\cdot\vec r=(2,-1,3)\cdot(2,-2,3)=4+2+9=15\,\text{J}.
\]
\item Si $W=0$, entonces $\vec F\cdot\vec r=0$. Si ninguno de los vectores es nulo, las direcciones de la fuerza y del desplazamiento son perpendiculares.
\item La fuerza de $10\,\text{N}$ en la dirección de $\vec u=(3,4)$ es
\[
\vec F=10\frac{(3,4)}{\sqrt{3^2+4^2}}=10\frac{(3,4)}5=(6,8).
\]
El desplazamiento es
\[
\vec r=Q-P=(6,0)-(0,0)=(6,0).
\]
Entonces
\[
W=\vec F\cdot\vec r=(6,8)\cdot(6,0)=36\,\text{J}.
\]
\end{enumerate}
\end{sol}

\item Sea $\vec u=(5,-12)$.
\begin{sol}
\[
\|\vec u\|=\sqrt{5^2+(-12)^2}=13.
\]
\begin{enumerate}[a)]
\item Un vector perpendicular se obtiene intercambiando componentes y cambiando un signo:
\[
\vec v=(12,5).
\]
Verificamos:
\[
\vec u\cdot\vec v=5\cdot12+(-12)5=0,\qquad \|\vec v\|=\sqrt{12^2+5^2}=13.
\]
\item A partir del versor perpendicular $\dfrac{(12,5)}{13}$, un vector ortogonal de módulo $10$ es
\[
\vec w=10\frac{(12,5)}{13}=\left(\frac{120}{13},\frac{50}{13}\right).
\]
\end{enumerate}
\begin{center}
\begin{tikzpicture}[>=triangle 45,scale=.35]
\ejesXY{-1}{13}{-13}{6}
\dibvec[blue]{(0,0)}{(5,-12)}{$\vec u$}
\dibvec[red]{(0,0)}{(12,5)}{$\vec v$}
\dibvec[green!60!black]{(0,0)}{(9.23,3.85)}{$\vec w$}
\end{tikzpicture}
\end{center}
\end{sol}

\item Dados $\vec u=(1,0,2)$ y $\vec v=(2,-1,1)$.
\begin{sol}
\begin{enumerate}[a)]
\item Un vector perpendicular a ambos es el producto vectorial:
\[
\vec u\times\vec v=
\begin{vmatrix}
\iv&\jv&\kv\\
1&0&2\\
2&-1&1
\end{vmatrix}
=(2,3,-1).
\]
\item Todos los vectores perpendiculares a ambos son los paralelos a $(2,3,-1)$:
\[
\vec n=\lambda(2,3,-1),\qquad \lambda\in\mathbb R,\qquad \lambda\neq0.
\]
\item Como $\|(2,3,-1)\|=\sqrt{14}$, para que $\|\lambda(2,3,-1)\|=7$ necesitamos
\[
|\lambda|\sqrt{14}=7\Rightarrow |\lambda|=\frac7{\sqrt{14}}=\frac{\sqrt{14}}{2}.
\]
Entonces
\[
\vec n_1=\frac{\sqrt{14}}2(2,3,-1)=\left(\sqrt{14},\frac{3\sqrt{14}}2,-\frac{\sqrt{14}}2\right)
\]
y su opuesto también cumple. No es único: hay dos con módulo $7$.
\end{enumerate}
\end{sol}

\item Paralelismo entre $\vec v=(a-2,3,-6)$ y $\vec w=(2b-4,-6,12)$.
\begin{sol}
\textbf{Método 1.} Si son paralelos, existe $\lambda$ tal que $\vec v=\lambda\vec w$. De las componentes segunda y tercera:
\[
3=\lambda(-6)\Rightarrow \lambda=-\frac12,
\qquad
-6=\lambda(12)\Rightarrow \lambda=-\frac12.
\]
Luego, con la primera componente:
\[
a-2=-\frac12(2b-4)=-b+2
\Rightarrow a+b=4.
\]

\textbf{Método 2.} En $\mathbb R^3$, dos vectores son paralelos si su producto vectorial es nulo.
\[
\vec v\times\vec w=
\begin{vmatrix}
\iv&\jv&\kv\\
a-2&3&-6\\
2b-4&-6&12
\end{vmatrix}.
\]
La primera componente es $3\cdot12-(-6)(-6)=36-36=0$. Las restantes dan una condición equivalente:
\[
-\big((a-2)12-(-6)(2b-4)\big)=0
\Rightarrow -12(a-2)-6(2b-4)=0
\Rightarrow a+b=4.
\]
Por lo tanto,
\[
\boxed{\vec v\parallel\vec w \Longleftrightarrow a+b=4.}
\]
\end{sol}

\item Producto vectorial opuesto.
\begin{sol}
Por la propiedad anticonmutativa del producto vectorial,
\[
\vec v\times\vec u=-(\vec u\times\vec v)=-(2,-4,1)=(-2,4,-1).
\]
\end{sol}

\item Área del paralelogramo.
\begin{sol}
Tenemos
\[
\vec a=2\iv+x\jv-\kv=(2,x,-1),
\qquad
\vec b=\iv+2\jv=(1,2,0).
\]
El área es $\|\vec a\times\vec b\|$. Calculamos:
\[
\vec a\times\vec b=
\begin{vmatrix}
\iv&\jv&\kv\\
2&x&-1\\
1&2&0
\end{vmatrix}
=(2,1,4-x).
\]
Entonces
\[
\|\vec a\times\vec b\|=\sqrt{2^2+1^2+(4-x)^2}=\sqrt{5+(4-x)^2}.
\]
Pedimos que el área sea $3$:
\[
\sqrt{5+(4-x)^2}=3
\Rightarrow 5+(4-x)^2=9
\Rightarrow (4-x)^2=4.
\]
Por lo tanto,
\[
4-x=2\Rightarrow x=2,
\qquad
4-x=-2\Rightarrow x=6.
\]
\[
\boxed{x=2\text{ o }x=6.}
\]
\end{sol}

\item Superficie total del cuerpo.
\begin{sol}
El cuerpo es un prisma triangular recto. La base triangular tiene catetos $2a$ y $3a$, y la longitud del prisma es $5a$.

\begin{center}
\begin{tikzpicture}[line join=round,line cap=round,scale=.65,x={(.9cm,-.25cm)},y={(.55cm,.35cm)},z={(0cm,.85cm)},>=triangle 45]
\coordinate (A) at (0,0,0);\coordinate (B) at (5,0,0);\coordinate (C) at (0,2,0);\coordinate (D) at (5,2,0);
\coordinate (E) at (0,0,3);\coordinate (F) at (5,0,3);
\draw[->] (0,0,0)--(6,0,0) node[right]{$x$};
\draw[->] (0,0,0)--(0,3,0) node[above]{$y$};
\draw[->] (0,0,0)--(0,0,4) node[above]{$z$};
\draw[thick] (A)--(B)--(D)--(C)--cycle;
\draw[thick] (A)--(E)--(F)--(B);
\draw[thick] (C)--(E);
\draw[thick] (D)--(F);
\node[below] at (B){$5a$};\node[left] at (C){$2a$};\node[left] at (E){$3a$};
\end{tikzpicture}
\end{center}

Área de las dos bases triangulares:
\[
2\cdot\frac12(2a)(3a)=6a^2.
\]
Las tres caras laterales son rectángulos de lados:
\[
5a\cdot2a,\qquad 5a\cdot3a,\qquad 5a\cdot \sqrt{(2a)^2+(3a)^2}=5a\cdot a\sqrt{13}.
\]
Entonces
\[
S_T=6a^2+10a^2+15a^2+5\sqrt{13}a^2=(31+5\sqrt{13})a^2.
\]
\end{sol}

\item Coplanaridad.
\begin{sol}
\begin{enumerate}[a)]
\item Tomamos los vectores con origen común en $A$:
\[
\overrightarrow{AB}=B-A=(1,2,-1),
\]
\[
\overrightarrow{AC}=C-A=(2,-1,2),
\qquad
\overrightarrow{AD}=D-A=(4,3,-1).
\]
Calculamos
\[
\overrightarrow{AC}\times\overrightarrow{AD}=
\begin{vmatrix}
\iv&\jv&\kv\\
2&-1&2\\
4&3&-1
\end{vmatrix}=(-5,10,10).
\]
Luego,
\[
\overrightarrow{AB}\cdot(\overrightarrow{AC}\times\overrightarrow{AD})=(1,2,-1)\cdot(-5,10,10)=-5+20-10=5.
\]
Como el producto mixto es distinto de cero, los puntos no son coplanares.
\item Los vectores son
\[
\vec u=(1,-2,3),\qquad \vec v=(\alpha,1,-1),\qquad \vec w=(2,0,1).
\]
Deben cumplir $\vec u\cdot(\vec v\times\vec w)=0$.
\[
\vec v\times\vec w=
\begin{vmatrix}
\iv&\jv&\kv\\
\alpha&1&-1\\
2&0&1
\end{vmatrix}=(1,-\alpha-2,-2).
\]
Entonces
\[
\vec u\cdot(\vec v\times\vec w)=(1,-2,3)\cdot(1,-\alpha-2,-2)=1+2\alpha+4-6=2\alpha-1.
\]
Pedimos
\[
2\alpha-1=0\Rightarrow \boxed{\alpha=\frac12}.
\]
\end{enumerate}
\end{sol}

\item Volumen, área de base, altura y vector altura.
\begin{sol}
Recordamos:
\[
V=|\vec u\cdot(\vec v\times\vec w)|,
\qquad
A=\|\vec v\times\vec w\|,
\qquad
h=\frac{V}{A},
\]
\[
\vec h=\operatorname{Proy}_{\vec v\times\vec w}\vec u
=\frac{\vec u\cdot(\vec v\times\vec w)}{\|\vec v\times\vec w\|^2}(\vec v\times\vec w).
\]

\begin{enumerate}[i)]
\item $\vec u=(1,2,3)$, $\vec v=(2,0,1)$ y $\vec w=(1,3,0)$.
\[
\vec v\times\vec w=
\begin{vmatrix}
\iv&\jv&\kv\\
2&0&1\\
1&3&0
\end{vmatrix}=(-3,1,6).
\]
Producto mixto:
\[
\vec u\cdot(\vec v\times\vec w)=(1,2,3)\cdot(-3,1,6)=-3+2+18=17.
\]
Por lo tanto,
\[
V=17,\qquad A=\sqrt{(-3)^2+1^2+6^2}=\sqrt{46},\qquad h=\frac{17}{\sqrt{46}}.
\]
El vector altura es
\[
\vec h=\frac{17}{46}(-3,1,6)=\left(-\frac{51}{46},\frac{17}{46},\frac{51}{23}\right).
\]

\begin{center}
\begin{tikzpicture}[line join=round,line cap=round,scale=.75,>=triangle 45]
\coordinate (O) at (0,0);\coordinate (V) at (3,0);\coordinate (W) at (1,1.4);\coordinate (VW) at (4,1.4);\coordinate (U) at (1.2,3.3);
\draw[fill=gray!10] (O)--(V)--(VW)--(W)--cycle;
\draw[->,thick,blue] (O)--(V) node[midway,below]{$\vec v$};
\draw[->,thick,red] (O)--(W) node[midway,left]{$\vec w$};
\draw[->,thick,green!60!black] (O)--(U) node[midway,left]{$\vec u$};
\draw[dashed] (U)--(2,1.4);
\draw[->,thick,purple] (2,1.4)--(U) node[midway,right]{$\vec h$};
\node at (2.4,.7){\scriptsize base};
\end{tikzpicture}
\end{center}

\item $\vec u=(2,-2,3)$, $\vec v=(1,0,2)$ y $\vec w=(0,2,1)$.
\[
\vec v\times\vec w=
\begin{vmatrix}
\iv&\jv&\kv\\
1&0&2\\
0&2&1
\end{vmatrix}=(-4,-1,2).
\]
Producto mixto:
\[
\vec u\cdot(\vec v\times\vec w)=(2,-2,3)\cdot(-4,-1,2)=-8+2+6=0.
\]
Por lo tanto,
\[
V=0,\qquad A=\sqrt{16+1+4}=\sqrt{21},\qquad h=0,\qquad \vec h=\vec 0.
\]
En este caso los tres vectores son coplanares, por eso no determinan volumen.
\end{enumerate}
\end{sol}

\end{enumerate}
\end{document}
